\chapter{La codificación de \Compact{}}
\label{ch:compact-relation}

La estructura de datos más densa del motor es la representación
en memoria de las aristas. Un grafo de 200~M aristas a $\sim$200
bytes cada una (el \code{Relation} completo con metadatos, strings,
etc.) ocuparía 40~GB --- por eso el motor guarda cada arista dos
veces: una como \code{Relation} rica, usada sólo mientras la
arista se está insertando, y de ahí en más como \Compact{} de 32
bytes, que vive en la lista de adyacencia el resto de la sesión.

\section{Layout}

\Compact{} es \code{\#[repr(C)]}, con los campos ordenados por
mayor a menor alineación para evitar padding interior. La
Figura~\ref{fig:compact-layout} muestra el layout byte a byte.

\begin{figure}[h]
\centering
\begin{tikzpicture}[
    byte/.style={draw, minimum width=0.42cm, minimum height=0.7cm,
                 font=\scriptsize, inner sep=1pt},
    label/.style={font=\scriptsize},
    x=0.42cm,
]
\foreach \i in {0,...,7} {
    \node[byte, fill=blue!12] at (\i, 0) (ts\i) {};
}
\node[label, above=0pt of ts3.north]{\code{timestamp}: \code{i64}};
\node[label, below=0pt of ts0.south]{offset 0};

\foreach \i in {0,...,7} {
    \node[byte, fill=orange!18] at (\i, -1) (mo\i) {};
}
\node[label, above=0pt of mo3.north]{\code{metadata\_offset}: \code{u64}};
\node[label, below=0pt of mo0.south]{offset 8};

\foreach \i in {0,...,3} {
    \node[byte, fill=green!18] at (\i, -2) (ssid\i) {};
}
\foreach \i [evaluate=\i as \j using int(\i+4)] in {0,...,3} {
    \node[byte, fill=green!28] at (\j, -2) (dsid\i) {};
}
\node[label, above=0pt of ssid1.north]{\code{source\_sid}};
\node[label, above=0pt of dsid1.north]{\code{dest\_sid}};
\node[label, below=0pt of ssid0.south]{offset 16};
\node[label, below=0pt of dsid0.south]{offset 20};

\foreach \i in {0,...,3} {
    \node[byte, fill=red!18] at (\i, -3) (ia\i) {};
}
\foreach \i [evaluate=\i as \j using int(\i+4)] in {0,...,1} {
    \node[byte, fill=purple!18] at (\j, -3) (dt\i) {};
}
\node[byte, fill=yellow!28] at (6, -3) (rt) {};
\node[byte, fill=gray!25] at (7, -3) (pad) {};
\node[label, above=0pt of ia1.north]{\code{ingested\_at\_delta}};
\node[label, above=0pt of dt0.north]{\code{dataset}};
\node[label, above=0pt of rt.north]{\code{rel}};
\node[label, above=0pt of pad.north]{pad};
\node[label, below=0pt of ia0.south]{offset 24};
\node[label, below=0pt of dt0.south]{offset 28};
\node[label, below=0pt of rt.south]{offset 30};
\node[label, below=0pt of pad.south]{offset 31};
\end{tikzpicture}
\caption{Layout de memoria de \Compact{} (32 bytes). El \code{u32}
\code{ingested\_at\_delta} (P3) recupera 4 de los 5 bytes de
padding final del layout pre-P3.}
\label{fig:compact-layout}
\end{figure}

Semánticamente:
\begin{itemize}
    \item \code{timestamp}: tiempo de evento de la arista, en
        segundos epoch.
    \item \code{metadata\_offset}: offset dentro del
        \code{MetadataStore} externo (Capítulo~\ref{ch:metadata});
        cero significa ``sin metadatos''.
    \item \code{source\_sid}, \code{dest\_sid}: manejadores de
        32~bits al \code{StringInterner}
        (Capítulo~\ref{ch:interner}).
    \item \code{ingested\_at\_delta}: segundos desde el epoch MVCC
        base (ver abajo); cero en aristas heredadas.
    \item \code{dataset\_tag}: manejador de 16~bits a la tabla
        \code{Arc<str>} de identificadores de dataset por grafo;
        cero significa ``sin dataset''.
    \item \code{rel\_type\_tag}: enum tag de 8~bits para el tipo
        de relación.
\end{itemize}

El \code{rel\_type\_tag} es deliberadamente \code{u8} para que el
matcher pueda hacer comparaciones SIMD sin ramas
(Capítulo~\ref{ch:simd}).

\section{Codificación MVCC}

\begin{invariant}[Tamaño MVCC]
\code{size\_of::<CompactRelation>() == 32}. Esto se asevera en
\code{tests/relation\_schema.rs}. Violarlo cuesta $\sim$20\% del
footprint del grafo en memoria, así que es una aserción cargada
de peso.
\end{invariant}

El delta MVCC se mide contra
\code{INGEST\_EPOCH\_BASE} = 1\,704\,067\,200 segundos (2024-01-01
UTC). Un \code{u32} cubre $\approx 2^{32}$ segundos, unos 136
años; el motor satura en \code{u32::MAX} para ingestas más allá
de $\sim$2160. Las aristas heredadas cargadas desde un archivo de
sesión pre-P3 resuelven al propio epoch base, que es un instante
definido aunque impreciso.

El Algoritmo~\ref{alg:encode-ingested-at} muestra el codificador.
La decodificación es el trivial \code{INGEST\_EPOCH\_BASE + delta
as i64}.

\begin{algorithm}
\caption{\textsc{Codificar-Ingested-At}}\label{alg:encode-ingested-at}
\KwIn{timestamp absoluto en segundos epoch, $t$}
\KwOut{delta \code{u32} a guardar en \Compact{}}
\If{$t \leq $ \code{INGEST\_EPOCH\_BASE}}{
    \Return 0\;
}
$\delta \gets t - $ \code{INGEST\_EPOCH\_BASE}\;
\If{$\delta > $ \code{u32::MAX}}{
    \Return \code{u32::MAX} \Comment*[r]{saturar}
}
\Return $\delta$ como \code{u32}\;
\end{algorithm}

\begin{complexity}
Tiempo y espacio constantes. El almacenamiento por arista es de
32 bytes, dando $\Thet{32 m}$ para la lista de aristas de un
grafo con $m$ aristas.
\end{complexity}

\section{¿Por qué no \code{\#[repr(C, packed)]}?}

Un layout empaquetado ahorraría otros 5 bytes eliminando el pad
final, llegando a 27 bytes por arista ($\sim$16\% adicional). El
motor no lo usa porque:

\begin{enumerate}
    \item Cada lectura de un campo empaquetado requiere una copia
        byte a byte hacia storage alineado para satisfacer el
        invariante de alineación de referencias en Rust. El
        matcher DFS lee \code{source\_sid} y \code{dest\_sid}
        millones de veces por hunt; ese overhead de copia
        dominaría cualquier ahorro de memoria.
    \item La alineación a 32~bytes coincide con la unidad de load
        AVX2, así que los barridos SIMD de la lista de aristas se
        mantienen amigables a la línea de caché.
\end{enumerate}

Por eso el pad final queda como espacio futuro para más
extensiones MVCC (p.~ej.\ un puntaje de confianza o un bit flag
de tag de fuente).

\begin{inthecode}
\filepath{graph\_hunter\_core/src/relation.rs:83} --- definición
del struct.\\
\filepath{graph\_hunter\_core/src/relation.rs:151} --- helper
\code{encode\_ingested\_at}.\\
\filepath{graph\_hunter\_core/tests/mvcc\_ingested\_at.rs} ---
tests de round-trip MVCC.\\
Tests de módulo en \filepath{graph\_hunter\_core/src/relation.rs}
--- invariante de 32 bytes.
\end{inthecode}
